\documentclass[11pt, a4paper]{article}
\usepackage[utf8]{inputenc}
\usepackage{amsmath, amssymb}
\usepackage{geometry}
\usepackage{graphicx}
\usepackage{fancyhdr}
\usepackage{titlesec}
\usepackage{enumitem}
\usepackage{xcolor}
\usepackage{mdframed}
\usepackage{parskip}
\usepackage{tikz}
\usepackage{listings}
\usepackage{booktabs}
\usetikzlibrary{calc}

\geometry{top=2.5cm, bottom=2.5cm, left=2.5cm, right=2.5cm}
\definecolor{unsam}{RGB}{87, 122, 160}
\definecolor{lightgray}{RGB}{245, 245, 245}
\definecolor{conceptbg}{RGB}{235, 245, 255}
\definecolor{warnbg}{RGB}{255, 248, 230}
\definecolor{warnline}{RGB}{200, 150, 0}
\definecolor{codegray}{RGB}{250,250,250}
\definecolor{codeblue}{RGB}{0,80,160}
\definecolor{codegreen}{RGB}{0,120,0}
\definecolor{codecomment}{RGB}{100,100,100}

\lstdefinestyle{pythonstyle}{
  language=Python, backgroundcolor=\color{codegray},
  basicstyle=\footnotesize\ttfamily, commentstyle=\color{codecomment},
  keywordstyle=\color{codeblue}\bfseries, stringstyle=\color{codegreen},
  numberstyle=\tiny\color{gray}, breaklines=true, captionpos=b,
  numbers=left, numbersep=5pt, showstringspaces=false, tabsize=4,
  frame=single, rulecolor=\color{gray!40}, xleftmargin=15pt, xrightmargin=5pt,
}

\pagestyle{fancy}
\fancyhf{}
\fancyhead[L]{\small \textcolor{unsam}{\textbf{An\'alisis y Procesamiento de Se\~nales}}}
\fancyhead[R]{\small \textcolor{unsam}{UNSAM --- Ciencia y Tecnolog\'ia}}
\fancyfoot[C]{\small \thepage}
\renewcommand{\headrulewidth}{0.6pt}
\renewcommand{\footrulewidth}{0pt}

\newcommand{\sechdr}[1]{%
  \noindent\colorbox{unsam}{%
    \parbox{\dimexpr\textwidth\relax}{%
      \vspace{4pt}\hspace{6pt}{\large\bfseries\color{white}#1}\vspace{4pt}%
    }%
  }%
}
\titleformat{\section}[block]{\normalsize}{}{0em}{\sechdr}
\titlespacing*{\section}{0pt}{1.8em}{0.9em}

\newmdenv[backgroundcolor=conceptbg,linecolor=unsam,linewidth=1.5pt,
  innerleftmargin=12pt,innerrightmargin=12pt,innertopmargin=10pt,
  innerbottommargin=10pt,skipabove=8pt,skipbelow=8pt]{conceptbox}
\newmdenv[backgroundcolor=lightgray,linecolor=gray!35,linewidth=0.8pt,
  innerleftmargin=14pt,innerrightmargin=14pt,innertopmargin=11pt,
  innerbottommargin=11pt,skipabove=7pt,skipbelow=7pt]{exbox}
\newmdenv[backgroundcolor=warnbg,linecolor=warnline,linewidth=1.2pt,
  innerleftmargin=12pt,innerrightmargin=12pt,innertopmargin=8pt,
  innerbottommargin=8pt,skipabove=6pt,skipbelow=6pt]{hintbox}

\begin{document}

% ===== PORTADA =====
\begin{center}
  {\LARGE\bfseries\color{unsam} Gu\'ia G5}\\[0.25cm]
  {\large\bfseries Ventanas para An\'alisis Espectral}\\[0.15cm]
  {\color{unsam}\rule{0.6\linewidth}{1.2pt}}
\end{center}
\vspace{0.3cm}

\begin{conceptbox}
\textbf{Ideas clave} \hfill \small\textit{Holton, Cap.~14}
\medskip

Toda se\~nal analizada con la DFT es de \textbf{duraci\'on finita}: equivale a multiplicar
la se\~nal original por una \textbf{ventana} $w_N[n]$:
$$\hat{x}[n] = x[n]\cdot w_N[n]
\;\;\Longrightarrow\;\;
\hat{X}(\omega) = \frac{1}{2\pi}\,X(\omega) * W_N(\omega)$$
El espectro resultante es la \emph{convoluci\'on} entre el espectro verdadero y la
transformada de la ventana, lo que introduce dos efectos adversos:

\begin{itemize}[noitemsep]
  \item \textbf{Desparramo} (\textit{spectral spread}): el ancho del l\'obulo principal de
    $W_N(\omega)$ limita la \emph{resoluci\'on espectral}, es decir, la capacidad de separar
    dos componentes pr\'oximas en frecuencia.
  \item \textbf{Filtraci\'on} (\textit{spectral leakage}): los l\'obulos laterales de
    $W_N(\omega)$ crean un ``piso espectral'' que puede enmascarar componentes de peque\~na
    amplitud alejadas de la componente dominante.
\end{itemize}

\smallskip
\textbf{Par\'ametros clave de una ventana} (Holton, Sec.~14.1.2):

\smallskip
\begin{center}
\begin{tabular}{lcccc}
\toprule
\textbf{Ventana} & $\Delta\omega$ (l\'ob.\ principal) & $\Delta A_{\max}$ & Decaída & Punto fuerte \\
\midrule
Rectangular        & $4\pi/N$   & $-13\,\text{dB}$ & $-6\,\text{dB/oct}$  & Máxima resolución \\
Hann               & $8\pi/N$   & $-32\,\text{dB}$ & $-18\,\text{dB/oct}$ & Compromiso general \\
Hamming            & $8\pi/N$   & $-43\,\text{dB}$ & $-6\,\text{dB/oct}$  & Primer lóbulo mín. \\
Blackman-Harris    & $12\pi/N$  & $-92\,\text{dB}$ & $-$                  & Señales muy débiles \\
Flattop            & $\approx18\pi/N$ & $-93\,\text{dB}$ & $-$ & Medición de amplitud \\
\bottomrule
\end{tabular}
\end{center}
\end{conceptbox}

% =====================================================================
\section{Periodicidad Impl\'icita y Desparramo Espectral}

\begin{exbox}
\textbf{Ejercicio 1.}
Responda las siguientes preguntas sobre los fundamentos del an\'alisis espectral con DFT.

\begin{enumerate}[label=\textbf{\alph*)}]

  \item \textbf{Periodicidad impl\'icita.}
  La DFT asume que la se\~nal analizada se repite peri\'odicamente.
  Explique, en sus propias palabras, por qu\'e esta suposici\'on introduce
  \emph{discontinuidades artificiales} cuando la frecuencia de la se\~nal
  no corresponde exactamente a un bin de la DFT.
  Ayuda: piense qu\'e ocurre al ``pegar'' el \'ultimo y el primer valor del bloque analizado.

  \item \textbf{Spread vs.\ leakage.}
  Holton (Sec.~14.1.3) distingue dos efectos adversos del enventanado.
  Explique la diferencia entre \emph{spectral spread} y \emph{spectral leakage}:
  \begin{itemize}[noitemsep]
    \item ?`Cu\'al de los dos limita la resoluci\'on entre dos tonos cercanos?
    \item ?`Cu\'al de los dos impide detectar una componente d\'ebil cerca de una fuerte?
    \item ?`De qu\'e par\'ametro de la ventana depende cada uno?
  \end{itemize}

  \item \textbf{Ejemplo num\'erico.}
  Se tiene $x[n] = \cos(2\pi \cdot 205 \cdot n/f_s)$ con $f_s = 1000\,\text{Hz}$
  y se analizan $N = 100$ muestras.
  \begin{itemize}[noitemsep]
    \item Calcule el espaciado de bins: $\Delta f = f_s/N$.
    \item Determine si $f_0 = 205\,\text{Hz}$ cae exactamente en un bin.
      Exprese el desvío en fracci\'on de bin ($\Delta \in [-0.5,\, 0.5]$).
    \item Si se usara una ventana rectangular, ?`en qu\'e regi\'on del espectro
      aparecería energía que no pertenece a $f_0$?
      Describa cualitativamente usando los conceptos de spread y leakage.
  \end{itemize}

\end{enumerate}
\end{exbox}

% =====================================================================
\section{Propiedades y Elecci\'on de Ventanas}

\begin{exbox}
\textbf{Ejercicio 2.}
Utilice la tabla de la caja de ideas clave (Holton, Fig.~14.6) para responder.

\begin{enumerate}[label=\textbf{\alph*)}]

  \item \textbf{Resoluci\'on m\'inima.}
  La resoluci\'on espectral en Hz puede estimarse como:
  $$\Delta f_{\min} = \frac{\Delta\omega}{2\pi} \cdot f_s$$
  donde $\Delta\omega$ es el ancho del l\'obulo principal de la ventana.
  Para $f_s = 8000\,\text{Hz}$ y $N = 800$ muestras, calcule $\Delta f_{\min}$
  para cada una de las cinco ventanas de la tabla.
  ?`Cu\'al permite separar dos tonos a $50\,\text{Hz}$ de distancia?

  \item \textbf{Piso de leakage.}
  Suponga una se\~nal formada por dos sinusoidales: la primera de $0\,\text{dB}$
  y la segunda de $-40\,\text{dB}$, separadas $50\,\text{Hz}$.
  Para cada ventana, indique si los l\'obulos laterales de la se\~nal dominante
  ($0\,\text{dB}$) \emph{enmascaran} o \emph{no enmascaran} la se\~nal d\'ebil ($-40\,\text{dB}$).
  Use directamente el valor $\Delta A_{\max}$ de la tabla.

  \item \textbf{Ventana Flattop.}
  La ventana Flattop tiene el l\'obulo principal m\'as ancho de la tabla.
  Sin embargo, la comunidad de metrolog\'ia la usa ampliamente.
  ?`Para qu\'e aplicaci\'on resulta \'util? ?`Qu\'e propiedad del l\'obulo principal
  justifica su uso?

  \item \textbf{Decaída de lóbulos laterales.}
  La columna ``decaída'' indica cu\'an r\'apido caen los l\'obulos laterales
  a medida que nos alejamos de la frecuencia de la se\~nal.
  Compare Hamming ($-6\,\text{dB/oct}$) vs.\ Hann ($-18\,\text{dB/oct}$):
  ambas tienen el mismo $\Delta A_{\max}$ del \emph{primer} l\'obulo lateral,
  pero se comportan diferente a frecuencias distantes.
  ?`Cu\'al conviene si la se\~nal d\'ebil est\'a \emph{lejos} en frecuencia de la dominante?

\end{enumerate}
\end{exbox}

% =====================================================================
\section{Sego en Estimaci\'on de Amplitud: Ventana Rectangular vs.\ Flattop}

\begin{conceptbox}
\textbf{Ganancia coherente (CG) y correci\'on de amplitud} \hfill \small\textit{Holton, Sec.~14.1.2}

\smallskip
Dado un espectro ventaneado $X_w[k] = \text{FFT}\{x[n]\cdot w[n]\}$, el pico espectral en
la frecuencia de una sinusoide de amplitud $A$ vale, aproximadamente:
$$|X_w[k_0]| = \frac{N}{2}\cdot \text{CG}\cdot A \cdot |\text{sinc}(\Delta)|$$
donde la \textbf{ganancia coherente} $\text{CG} = \tfrac{1}{N}\sum_{n=0}^{N-1}w[n]$
y $\Delta \in [-0.5,\,0.5]$ es el desvío en fracci\'on de bin entre $f_0$ y el bin m\'as
cercano.

Para recuperar la amplitud verdadera $A$ se aplica la correcci\'on:
$$\hat{A} = \frac{2\,|X_w[k_0]|}{N\cdot\text{CG}}$$
Esta correcci\'on \emph{solo es exacta si $\Delta = 0$} (la frecuencia cae en un bin
exacto) o si la ventana tiene l\'obulo principal \emph{plano} (Flattop).
\end{conceptbox}

\begin{exbox}
\textbf{Ejercicio 3.}
Sea $x[n] = A\cdot\sin\!\bigl(\Omega_1\cdot n\bigr)$ con $A = 2$,
$\Omega_0 = \pi/2$ y $\Omega_1 = \Omega_0 + \Delta\cdot\tfrac{2\pi}{N}$,
donde $\Delta$ es un desvío fraccional ($\Delta = 0$: frecuencia exactamente en bin;
$\Delta = 0.5$: exactamente a mitad entre dos bins).

\begin{enumerate}[label=\textbf{\alph*)}]

  \item \textbf{Ventana rectangular, $\Delta = 0$.}
  Calcule el valor te\'orico de $|X_w[k_0]|$ para la ventana rectangular
  ($\text{CG}_{\text{rect}} = 1$) y $N = 1000$.
  Verifique que la correci\'on $\hat{A} = 2|X_w|/(N\cdot\text{CG})$ recupera $A = 2$.

  \item \textbf{Ventana rectangular, $\Delta = 0.5$.}
  Cuando la frecuencia cae a mitad de bin, el factor de atenuaci\'on es:
  $$|\text{sinc}(\Delta)| = \left|\frac{\sin(\pi\Delta)}{\pi\Delta}\right|
  \bigg|_{\Delta=0.5} = \frac{2}{\pi} \approx 0.637$$
  ?`Qu\'e amplitud $\hat{A}$ estima la ventana rectangular?
  Exprese el error en dB.
  ?`Es un error de \emph{sesgo} o de \emph{varianza}? ?`Por qu\'e?

  \item \textbf{Ventana Flattop.}
  La ventana Flattop se dise\~na para que $|\text{sinc}(\Delta)| \approx 1$
  para cualquier $|\Delta| \leq 0.5$.
  Explique \emph{cualitativamente} (sin calcular) c\'omo logra esto:
  ?`c\'omo debe ser la forma de su l\'obulo principal en frecuencia?
  ?`Qu\'e sacrifica para lograrlo?

  \item \textbf{Conexi\'on con la TS4.}
  En la TS4 se estudia el estimador de amplitud $\hat{a}_1 = |X_w^i(\Omega_0)|$
  para distintas ventanas con $\Omega_1 = \Omega_0 + f_r\cdot\tfrac{2\pi}{N}$,
  donde $f_r \sim \mathcal{U}(-2,\,2)$.
  \begin{itemize}[noitemsep]
    \item ?`Qu\'e ventana producir\'a menor \emph{sesgo} en la estimaci\'on de amplitud?
    \item ?`Qu\'e ventana producir\'a menor \emph{varianza} en la estimaci\'on de frecuencia?
    \item ?`Son las mismas? ?`Qu\'e trade-off refleja esto?
  \end{itemize}

\end{enumerate}
\end{exbox}

% =====================================================================
\section{Ejercicio Integrador: Dos Se\~nales Pr\'oximas}

\begin{hintbox}
\textbf{Contexto:}
Un escenario t\'ipico en instrumentaci\'on es la presencia de dos sinusoidales
con muy distinta amplitud y frecuencias pr\'oximas, sumergidas en ruido.
La elecci\'on de ventana, $N$ y $f_s$ determina si podemos
\emph{resolver} ambas componentes \emph{y} detectar la m\'as d\'ebil sin que la m\'as
fuerte la enmascare. Este tipo de problema aparece en los parciales de la materia.
\end{hintbox}

\begin{exbox}
\textbf{Ejercicio 4.}
Se tiene una se\~nal:
$$x(t) = \underbrace{A_1\cos(2\pi f_1 t)}_{0\,\text{dB},\; f_1=1000\,\text{Hz}}
       + \underbrace{A_2\cos(2\pi f_2 t)}_{-40\,\text{dB},\; f_2=1050\,\text{Hz}}
       + n(t)$$
donde el ruido $n(t)$ es gaussiano con nivel de piso $-60\,\text{dB}$.
El sistema de adquisici\'on trabaja a $f_s = 8000\,\text{Hz}$.

\begin{enumerate}[label=\textbf{\alph*)}]

  \item \textbf{Condici\'on de resoluci\'on.}
  La separaci\'on entre $f_1$ y $f_2$ es $\Delta f = 50\,\text{Hz}$.
  Para cada ventana de la tabla del Ejercicio~2, determine el $N$ m\'inimo
  tal que $\Delta f_{\min} < \Delta f$.
  Organice los resultados en una tabla de dos columnas: ventana y $N_{\min}$.

  \item \textbf{Condici\'on de leakage.}
  Independientemente de $N$, analice si cada ventana puede ``ver'' $A_2$ por encima
  del piso de l\'obulos laterales de $A_1$ (que est\'a $40\,\text{dB}$ por encima).
  Use la columna $\Delta A_{\max}$ directamente.
  Indique para cada ventana: ¿resuelve? ¿detecta $A_2$?

  \item \textbf{Ventana \'optima.}
  Seleccione la ventana que satisfaga \emph{ambas} condiciones
  (resoluci\'on y deteci\'on de $A_2$) con el menor $N$ posible.
  Proponga un valor de $N$ pr\'actico (potencia de 2, mayor que $N_{\min}$)
  y calcule la resoluci\'on efectiva $\Delta f = f_s / N$ que obtendr\'ia.

  \item \textbf{Pregunta conceptual de parcial.}
  Suponga ahora que el nivel de ruido sube de $-60\,\text{dB}$ a $-30\,\text{dB}$.
  ?`Alcanza con cambiar la ventana? ?`Qu\'e otras alternativas (relacionadas con $N$
  o con el m\'etodo de estimaci\'on) podr\'ia explorar?
  Mencione al menos dos estrategias y c\'omo afectan sesgo y varianza del estimador.

\end{enumerate}
\end{exbox}

% =====================================================================
\section{Python: An\'alisis de Nota Musical}

\begin{exbox}
\textbf{Ejercicio 5.} (Exploratorio con c\'odigo)

Se proveen los archivos \texttt{La440.wav} (nota La a 440~Hz) y
\texttt{La440\_delay.wav} (misma nota con retardo).
Usando el c\'odigo de la secci\'on siguiente como base:

\begin{enumerate}[label=\textbf{\alph*)}]

  \item \textbf{Respuesta en frecuencia de las ventanas.}
  Habilite las cinco ventanas del c\'odigo (descomentando las l\'ineas correspondientes)
  y grafique $|W_N(\omega)|_{\text{dB}}$ superpuestas.
  Identifique visualmente el ancho del l\'obulo principal y el nivel del primer
  l\'obulo lateral para cada una. Compare con la tabla te\'orica.

  \item \textbf{Estimaci\'on de frecuencia fundamental.}
  Para cada ventana, el c\'odigo imprime la frecuencia del pico.
  ?`Todas coinciden en $\approx 440\,\text{Hz}$?
  ?`Cu\'al tiene mayor dispersi\'on visualmente en el espectro alrededor del pico?

  \item \textbf{Estimaci\'on de amplitud con correcci\'on CG.}
  El c\'odigo ya aplica la correcci\'on $X / (N \cdot \text{CG})$.
  Compare la amplitud estimada del pico entre ventanas.
  ?`Cu\'al da el valor m\'as estable frente al hecho de que 440~Hz
  puede \emph{no} caer exactamente en un bin seg\'un el $N$ disponible?
  Ayuda: calcule el desvío $\Delta$ para la $f_s$ y $N$ del archivo.

  \item \textbf{(Opcional) Comparaci\'on con \texttt{La440\_delay.wav}.}
  Cargue ambos archivos y superponga sus espectros (m\'odulo) usando la misma ventana.
  ?`Se puede distinguir el retardo a partir del espectro de m\'odulo?
  ?`Qu\'e informaci\'on del espectro s\'i reflejar\'ia el retardo?

\end{enumerate}
\end{exbox}

% =====================================================================
\newpage
\section{C\'odigo Python: An\'alisis con Ventanas}

\begin{hintbox}
\textbf{Para el Ejercicio~5:} el par\'ametro \texttt{sym=False} genera una ventana
\emph{peri\'odica} (apropiada para an\'alisis espectral); \texttt{sym=True} genera
una ventana \emph{sim\'etrica} (usada en dise\~no de filtros FIR).
No los intercambie.
\end{hintbox}

\begin{lstlisting}[style=pythonstyle, caption={ventanas\_analisis.py}]
import numpy as np
import matplotlib.pyplot as plt
from scipy.io import wavfile
import scipy.signal.windows as spsw

# ===== Leer WAV =====
fs, data = wavfile.read("La440.wav")
if data.ndim == 2:          # stereo -> mono
    data = data.mean(axis=1)
x = data.astype(np.float64)
x = x - x.mean()           # quitar componente DC

N   = len(x)
Ts  = 1.0 / fs

# ===== Definicion de ventanas (sym=False: periodica) =====
ventanas = {
    "Rectangular":      np.ones(N),
    "Hann":             spsw.hann(N, sym=False),
    "Hamming":          spsw.hamming(N, sym=False),
    "Blackman-Harris":  spsw.blackmanharris(N, sym=False),
    "Flattop":          spsw.flattop(N, sym=False),
}

# ===== Eje de frecuencias unilateral =====
freqs = np.fft.rfftfreq(N, d=Ts)   # [0, fs/2] en Hz

# ===== Respuesta en frecuencia de las ventanas (zero-padding) =====
zp        = 16 * N
freqs_win = np.fft.rfftfreq(zp, d=Ts)

plt.figure(figsize=(10, 5))
for nombre, w in ventanas.items():
    W     = np.fft.rfft(w, n=zp)
    W_db  = 20 * np.log10(np.abs(W) / (np.abs(W[0]) + 1e-14))
    plt.plot(freqs_win, W_db, label=nombre, lw=1.2)

plt.xlim(0, 1500)
plt.ylim(-120, 3)
plt.xlabel("Frecuencia [Hz]")
plt.ylabel("|W(f)| [dB] (normalizado a 0 dB)")
plt.title("Respuesta en frecuencia de las ventanas")
plt.grid(True, which="both")
plt.legend()
plt.tight_layout()

# ===== Espectros con correccion CG =====
for nombre, w in ventanas.items():
    CG  = w.mean()                     # Coherent Gain = (1/N) * sum(w)
    xw  = x * w
    X   = np.fft.rfft(xw) / (N * CG)  # correccion de amplitud
    mag = np.abs(X)

    f0_est = freqs[np.argmax(mag)]

    plt.figure(figsize=(9, 4))
    plt.plot(freqs, mag, lw=1, label=f"{nombre}  |  pico ~ {f0_est:.2f} Hz")
    plt.xlim(0, 1000)
    plt.xlabel("Frecuencia [Hz]")
    plt.ylabel("|X(f)| corregido (u.a.)")
    plt.title(f"Espectro con ventana {nombre}")
    plt.grid(True)
    plt.legend()
    plt.tight_layout()

# ===== ENBW: Equivalent Noise Bandwidth (en bins) =====
print(f"\n{'Ventana':<20} {'CG':>8} {'ENBW (bins)':>14}")
print("-" * 44)
for nombre, w in ventanas.items():
    CG   = w.mean()
    ENBW = N * np.sum(w**2) / np.sum(w)**2
    print(f"{nombre:<20} {CG:>8.4f} {ENBW:>14.4f}")

plt.show()
\end{lstlisting}

\begin{hintbox}
\textbf{Interpretaci\'on del ENBW:}
El \textbf{Equivalent Noise Bandwidth} (en bins) indica cu\'antos bins de la DFT
``ve'' efectivamente la ventana al estimar la potencia de ruido.
Un ENBW mayor implica que la estimaci\'on de potencia incluye m\'as ruido,
lo que aumenta la varianza del estimador.
Para la ventana rectangular, $\text{ENBW} = 1\,\text{bin}$;
para Hann, $\text{ENBW} \approx 1.5\,\text{bins}$.
El ENBW en Hz se obtiene como $\text{ENBW}_{\text{Hz}} = \text{ENBW}_{\text{bins}} \cdot (f_s / N)$.
\end{hintbox}

\end{document}
